\begin{examplebox}
\textbf{\textcolor{CExample}{例1（基础）}}

\textbf{\textcolor{CExample}{题目：}} 在等比数列 $\{a_n\}$ 中，已知 $a_2=3$，$a_8=12$，求 $a_4\cdot a_6$ 的值。

\textbf{\textcolor{CExample}{解题步骤：}}
\begin{description}[style=nextline, leftmargin=2.8em, labelsep=0.5em, itemsep=0.45em]
    \item[\textbf{第一步（找下标和）：}] 计算 $2+8=10$，$4+6=10$，所以 $2+8=4+6$。
    \par\textbf{理由：} 验证下标和相等。
    \item[\textbf{第二步（套用结论）：}] 由结论 $a_2\cdot a_8 = a_4\cdot a_6$。
    \par\textbf{理由：} 下标和相等时乘积相等。
    \item[\textbf{第三步（代入计算）：}] $a_4\cdot a_6 = 3\times 12 = 36$。
    \par\textbf{理由：} 直接乘法。
\end{description}

\textbf{\textcolor{CExample}{关键结论：}} \textbf{$a_4\cdot a_6 = 36$}。

\medskip
\textbf{\textcolor{CExample}{例2（稍变形）}}

\textbf{\textcolor{CExample}{题目：}} 在等比数列 $\{a_n\}$ 中，$a_3\cdot a_7 = 25$，求 $a_5$ 的值。

\textbf{\textcolor{CExample}{解题步骤：}}
\begin{description}[style=nextline, leftmargin=2.8em, labelsep=0.5em, itemsep=0.45em]
    \item[\textbf{第一步（找下标关系）：}] $3+7=10$，$5+5=10$，所以 $3+7=5+5$。
    \par\textbf{理由：} 下标和相等。
    \item[\textbf{第二步（套用结论）：}] 由结论 $a_3\cdot a_7 = a_5\cdot a_5 = a_5^2$。
    \par\textbf{理由：} 推广结论中 $m=n=5$ 时 $a_5\cdot a_5 = a_5^2$。
    \item[\textbf{第三步（解方程）：}] $a_5^2 = 25$，所以 $a_5 = \pm 5$。
    \par\textbf{理由：} 开平方需考虑正负，公比可正可负。
\end{description}

\textbf{\textcolor{CExample}{关键结论：}} \textbf{$a_5 = \pm 5$}。

\medskip
\textbf{\textcolor{CExample}{例3（综合应用）}}

\textbf{\textcolor{CExample}{题目：}} 在等比数列 $\{a_n\}$ 中，$a_1\cdot a_9 = 64$，且 $a_3 + a_7 = 20$，求公比 $q$ 的值。

\textbf{\textcolor{CExample}{解题步骤：}}
\begin{description}[style=nextline, leftmargin=2.8em, labelsep=0.5em, itemsep=0.45em]
    \item[\textbf{第一步（用结论求中间项）：}] 由 $1+9=10$，$5+5=10$，得 $a_1\cdot a_9 = a_5^2 = 64$，所以 $a_5 = \pm 8$。
    \par\textbf{理由：} 推广结论。
    \item[\textbf{第二步（表示另外两项）：}] $a_3 = a_5 / q^2$，$a_7 = a_5 \cdot q^2$。
    \par\textbf{理由：} 等比数列项间关系：$a_5 = a_3 q^2$，$a_7 = a_5 q^2$。
    \item[\textbf{第三步（代入求和）：}] $a_3 + a_7 = \frac{a_5}{q^2} + a_5 q^2 = a_5\left(\frac{1}{q^2}+q^2\right) = 20$。
    \par\textbf{理由：} 代入表达式。
    \item[\textbf{第四步（分类讨论）：}] 若 $a_5=8$，则 $\frac{1}{q^2}+q^2 = \frac{20}{8}=\frac{5}{2}$，解得 $q^2=2$ 或 $q^2=\frac{1}{2}$，故 $q=\pm\sqrt{2}$ 或 $q=\pm\frac{\sqrt{2}}{2}$；若 $a_5=-8$，则 $\frac{1}{q^2}+q^2 = \frac{20}{-8} = -\frac{5}{2}$，方程 $t+\frac{1}{t}=-\frac{5}{2}$ 无正实数解，故 $a_5=-8$ 舍去。
    \par\textbf{理由：} 解方程并验证可行性。
\end{description}

\textbf{\textcolor{CExample}{关键结论：}} \textbf{$q = \pm\sqrt{2}$ 或 $q = \pm\frac{\sqrt{2}}{2}$}。
\end{examplebox}
